
\documentclass{article}
\usepackage{colortbl}
\usepackage{makecell}
\usepackage{multirow}
\usepackage{supertabular}

\begin{document}

\newcounter{utterance}

\centering \large Interaction Transcript for game `eqbench', experiment `eq\_bench', episode 102 with reflexion{-}llama{-}3.1{-}8b{-}instruct.
\vspace{24pt}

{ \footnotesize  \setcounter{utterance}{1}
\setlength{\tabcolsep}{0pt}
\begin{supertabular}{c@{$\;$}|p{.15\linewidth}@{}p{.15\linewidth}p{.15\linewidth}p{.15\linewidth}p{.15\linewidth}p{.15\linewidth}}
    \# & \multicolumn{2}{c}{Player} && \multicolumn{2}{c}{Game Master} \\
    \hline

    \theutterance \stepcounter{utterance}  
    & & & \multicolumn{4}{p{0.6\linewidth}}{
        \cellcolor[rgb]{0.9,0.9,0.9}{
            \makecell[{{p{\linewidth}}}]{
                \texttt{\tiny{[P1$\langle$GM]}}
                \\ 
\texttt{Your task is to predict the likely emotional responses of a character in this dialogue:} \\
\\ 
\texttt{Elijah: So, you just decided to kill off Amelia, just like that?} \\
\texttt{Jefferson: It wasn't "just like that," Elijah. It was a carefully considered decision.} \\
\texttt{Elijah: Carefully considered? She was the heart of your story. She was the reason I fell in love with your books!} \\
\texttt{Jefferson: And she still is. Her death doesn't erase her existence.} \\
\texttt{Elijah: But it erases her future! You took that away from her, from us.} \\
\texttt{Jefferson: I didn't take anything away, Elijah. It's a story. It's fiction.} \\
\texttt{Elijah: Fiction impacts reality, Jefferson. You of all people should know that.} \\
\texttt{Jefferson: And I do. But it's my story to tell. My reality to create.} \\
\texttt{Elijah: Your reality? What about the readers? Don't we have a say?} \\
\texttt{Jefferson: A say? No. An interpretation? Yes.} \\
\texttt{Elijah: Interpretation? Is that what you call it when you rip our hearts out?} \\
\texttt{Jefferson: It's not about ripping hearts out. It's about evoking emotion.} \\
\texttt{Elijah: Well, congratulations, you've evoked plenty of emotion. Just not the kind you'd want.} \\
\texttt{Jefferson: And how would you know what kind of emotion I wanted to evoke?} \\
\texttt{Elijah: Because I know your work, Jefferson. I've read every single one of your books, cover to cover. I've lived in your worlds, with your characters. And this...this isn't you.} \\
\texttt{Jefferson: Maybe it's a side of me you've never seen before.} \\
\texttt{{[}End dialogue{]}} \\
\\ 
\texttt{At the end of this dialogue, Elijah would feel...} \\
\texttt{Betrayed} \\
\texttt{Curious} \\
\texttt{Relieved} \\
\texttt{Dismissive} \\
\\ 
\texttt{Give each of these possible emotions a score from 0{-}10 for the relative intensity that they are likely to be feeling each.} \\
\\ 
\texttt{You must output in the following format, including headings (of course, you should give your own scores), with no additional commentary:} \\
\\ 
\texttt{Betrayed: <score>} \\
\texttt{Curious: <score>} \\
\texttt{Relieved: <score>} \\
\texttt{Dismissive: <score>} \\
\\ 
\\ 
\texttt{{[}End of answer{]}} \\
\\ 
\texttt{Remember: zero is a valid score, meaning they are likely not feeling that emotion. You must score at least one emotion > 0.} \\
\\ 
\texttt{Your answer:} \\
            }
        }
    }
     \\ \\

    \theutterance \stepcounter{utterance}  
    & \multicolumn{4}{p{0.6\linewidth}}{
        \cellcolor[rgb]{0.561,0.737,0.561}{
            \makecell[{{p{\linewidth}}}]{
                \texttt{\tiny{[P1$\rangle$GM]}}
                \texttt{Betrayed: 9} \\
\texttt{Curious: 2} \\
\texttt{Relieved: 0} \\
\texttt{Dismissive: 0} \\
            }
        }
    }
    & & \\ \\

    \theutterance \stepcounter{utterance}  
    & & & \multicolumn{2}{p{0.3\linewidth}}{
        \cellcolor[rgb]{0.9,0.9,0.9}{
            \makecell[{{p{\linewidth}}}]{
                \texttt{\tiny{[GM$|$GM]}}
                \texttt{Betrayed: 9} \\
\texttt{Curious: 2} \\
\texttt{Relieved: 0} \\
\texttt{Dismissive: 0} \\
            }
        }
    }
    & & \\ \\

    \theutterance \stepcounter{utterance}  
    & & & \multicolumn{2}{p{0.3\linewidth}}{
        \cellcolor[rgb]{0.9,0.9,0.9}{
            \makecell[{{p{\linewidth}}}]{
                \texttt{\tiny{[GM$|$GM]}}
                \texttt{Betrayed: 8} \\
\texttt{Curious: 4} \\
\texttt{Relieved: 0} \\
\texttt{Dismissive: 2} \\
            }
        }
    }
    & & \\ \\

    \theutterance \stepcounter{utterance}  
    & & & \multicolumn{2}{p{0.3\linewidth}}{
        \cellcolor[rgb]{0.9,0.9,0.9}{
            \makecell[{{p{\linewidth}}}]{
                \texttt{\tiny{[GM$|$GM]}}
                \texttt{game\_result = LOSE} \\
            }
        }
    }
    & & \\ \\

    \theutterance \stepcounter{utterance}  
    & & \multicolumn{4}{p{0.6\linewidth}}{
        \cellcolor[rgb]{0.95,0.95,0.95}{
            \makecell[{{p{\linewidth}}}]{
                \texttt{\tiny{[P1$\leftrightarrow$P]}}
                \texttt{Post{-}episode Reflexion (not sent to GM):} \\
\texttt{New plan: } \\
\\ 
\texttt{Upon re{-}reading the dialogue, I realized that I focused too much on the surface{-}level emotions expressed by Elijah and didn't consider the underlying emotional dynamics. I should have paid more attention to the power struggle between Elijah and Jefferson, as well as Elijah's emotional investment in the story. } \\
\\ 
\texttt{To improve my analysis, I will consider the following:} \\
\\ 
\texttt{1. Identify the emotional stakes for Elijah: What does he stand to lose or gain from Jefferson's decision? How does this impact his emotional response?} \\
\texttt{2. Analyze the power dynamics: Who has control over the narrative, and how does this influence Elijah's perception of the situation?} \\
\texttt{3. Examine the emotional language: What words and phrases does Elijah use to describe his emotions, and how do these words reflect his emotional state?} \\
\texttt{4. Consider the context: What is the significance of Amelia's death in the story, and how does this impact Elijah's emotional response?} \\
\\ 
\texttt{By taking a more nuanced approach, I will be able to provide a more accurate analysis of Elijah's emotional state.} \\
            }
        }
    }
    & \\ \\

\end{supertabular}
}

\end{document}
